\documentclass{article}

\usepackage{makecell}

\title{Ch. 15 CW \#1}
\author{Principles of Microeconomics}
\date{\today}

\begin{document}

\maketitle

\section*{Instructions}

Please answer these questions on another sheet of paper and turn them in on Gradescope by next class.

\section*{Warm Up}

\begin{enumerate}

\item Throwback to Ch. 4: What two characteristics make a market competitive?

\item Throwback to Ch. 14: What was our core assumption about firms? What do we assume they're trying to accomplish?

\end{enumerate}

\section*{15.1. What Is a Competitive Market?}

\begin{enumerate}

	\item A perfectly competitive firm
	
		\begin{enumerate}
		
		\item chooses its price to maximize profits.
		
		\item sets its price to undercut other firms selling similar products.
		
		\item takes its price as given by market conditions.
		
		\item picks the price that yields the largest market share.
		
		\end{enumerate}
	
	\item Minnesota Milk is a thriving dairy company. It sells milk by the gallon, and the market price of a gallon is \$6.
	
		\begin{enumerate}
		
		\item Why is it reasonable to assume that the milk market is competitive?
	
		\item Please fill in the following table:
	
			\begin{center}
			\begin{tabular}{c | c | c | c | c }
			Quantity & Price & \makecell{Total\\Revenue} & \makecell{Average\\Revenue} & \makecell{Marginal\\Revenue} \\
			\hline
			1 gallon & \$6 & & & \\
			\hline
			2 & 6 & & & \\
			\hline
			3 & 6 & & & \\
			\hline
			4 & 6 & & & \\
			\hline
			5 & 6 & & & \\
			\hline
			6 & 6 & & & \\
			\hline
			7 & 6 & & & \\
			\hline
			8 & 6 & & &
			\end{tabular}
			\end{center}
			
		\end{enumerate}
		
	\item When a perfectly competitive firm increases the quantity it produces and sells by 10 percent, its marginal revenue \underline{\phantom{xxxxxxxxxx}} and its total revenue rises by \underline{\phantom{xxxxxxxxxx}}.
	
		\begin{enumerate}
		
		\item falls; less than 10 percent
		
		\item falls; exactly 10 percent
		
		\item stays the same; less than 10 percent
		
		\item stays the same; exactly 10 percent
		
		\end{enumerate}

\end{enumerate}

\section*{15.2. Profit Maximization and the Competitive Firm's Supply Curve}

\begin{enumerate}

\item Minnesota Milk is still thriving, and the market price of a gallon of milk is still \$6. 

	\begin{enumerate}
	
	\item Please fill in the following table.
	
		\begin{center}
		\begin{tabular}{c | c | c | c | c | c | c}
		Quantity & \makecell{Total\\Revenue} & \makecell{Total\\Cost} & Profit & \makecell{Marginal\\Revenue} & \makecell{Marginal\\Cost} & \makecell{Change\\in Profit} \\
		\hline
		0 gallons & \$0 & \$3 & & & & \\
		\hline
		1 & 6 & 5 & & & & \\
		\hline
		2 & 12 & 8 & & & & \\
		\hline
		3 & 18 & 12 & & & & \\
		\hline
		4 & 24 & 17 & & & & \\
		\hline
		5 & 30 & 23 & & & & \\
		\hline
		6 & 36 & 30 & & & & \\
		\hline
		7 & 42 & 38 & & & & \\
		\hline
		8 & 48 & 47 & & & &
		\end{tabular}
		\end{center}
		
	\item How many gallons should Minnesota Milk sell? There are two possible answers.
	
	\item What is true of marginal revenue and marginal cost at the profit maximizing quantities?
	
	\item What is the relationship between marginal revenue and marginal cost when change in profit is positive? What about when it's negative?
	
	\item Keeping in mind your answer to (c), try to explain your answer to part (b).
	
	\end{enumerate}
	
\item A competitive firm maximizes profit by choosing the quantity at which

	\begin{enumerate}
	
	\item average total cost is at its minimum.
	
	\item marginal cost equals the price.
	
	\item average total cost equals the price.
	
	\item marginal cost equals average total cost.
	
	\end{enumerate}

\end{enumerate}

\end{document}